\chapter{Tail PBP, Prop 11.3--11.4}
\label{chap:tail}

Reference: [BMSZb] Section~10.5 (tail definition); [BMSZ] Prop~11.3, 11.4;
internal blueprint \texttt{tail\_definition.md}, \texttt{lemma\_11\_5\_PBP\_identity.md}.

\section{The tail PBP}

\begin{definition}[Tail length $k$]
  \label{def:tailLen}
  \lean{PBP.tailLen_D, PBP.tailLen_B}
  \leanok
  \uses{def:PBP}
  For $\tau \in \PBP_D$ with $\mathcal{P}$ col~0 longer than $\mathcal{Q}$ col~0, the
  \emph{tail length} is
  \[ k := \tau.\mathcal{P}.\mathrm{colLen}(0) - \tau.\mathcal{Q}.\mathrm{colLen}(0). \]
  Equivalently $k = (\mathbf{r}_1(\check{\mathcal{O}}) - \mathbf{r}_2(\check{\mathcal{O}}))/2 + 1$ via the BV dual correspondence.
\end{definition}

\begin{definition}[Tail symbol $x_\tau$ and top-of-tail symbol $x_1$]
  \label{def:tailSymbol}
  \lean{PBP.tailSymbol_D, PBP.tailSymbol_B}
  \leanok
  \uses{def:tailLen}
  For D type, the tail is the sequence of cells
  $(\mathcal{Q}.\mathrm{colLen}(0), 0),\ (\mathcal{Q}.\mathrm{colLen}(0) + 1, 0),\ \ldots,\ (\mathcal{P}.\mathrm{colLen}(0) - 1, 0)$
  in $\mathcal{P}$'s column~0. We write:
  \begin{itemize}
    \item $x_\tau := \mathcal{P}_\tau(\mathcal{P}.\mathrm{colLen}(0) - 1, 0)$ — \emph{bottom} tail cell (maximum $\mathrm{layerOrd}$ by layer monotonicity).
    \item $x_1 := \mathcal{P}_\tau(\mathcal{Q}.\mathrm{colLen}(0), 0)$ — \emph{top} tail cell (minimum $\mathrm{layerOrd}$).
  \end{itemize}
  For B type: analogous symbols in $\mathcal{Q}$.
\end{definition}

\begin{definition}[Tail signature $\mathrm{tailSig}(\tau) = (p_{\tau_t}, q_{\tau_t})$]
  \label{def:tailSignature}
  \lean{PBP.tailSignature_D, PBP.tailSignature_B}
  \leanok
  \uses{def:tailLen, def:tailContrib}
  \[ \mathrm{tailSig}(\tau) := \sum_{i = \tau.\mathcal{Q}.\mathrm{colLen}(0)}^{\tau.\mathcal{P}.\mathrm{colLen}(0) - 1} \mathrm{tailContrib}(\mathcal{P}_\tau(i, 0)). \]
  Each cell contributes $\mathrm{tailContrib}$ (Def.~\ref{def:tailContrib}); components sum to~$2$ per cell.
\end{definition}

\begin{lemma}[Tail signature sum identity]
  \label{lem:tailSig_sum}
  \lean{PBP.tailSignature_D_sum_eq}
  \leanok
  \uses{def:tailSignature}
  $p_{\tau_t} + q_{\tau_t} = 2k$ where $k$ is the tail length.
\end{lemma}
\begin{proof}
  \leanok
  Each tail cell contributes $\mathrm{tailContrib}$ with $\mathrm{tailContrib}(\sigma).1 + \mathrm{tailContrib}(\sigma).2 = 2$ for every symbol $\sigma$. Summing over the $k$ tail cells gives $p_{\tau_t} + q_{\tau_t} = 2k$.
\end{proof}


\begin{lemma}[Tail signature non-negativity]
  \label{lem:tailSig_nonneg}
  \lean{tailSig_nonneg_D}
  \leanok
  \uses{def:tailSignature}
  $p_{\tau_t} \geq 0$ and $q_{\tau_t} \geq 0$.
\end{lemma}
\begin{proof}
  \leanok
  Each $\mathrm{tailContrib}(\sigma)$ has non-negative components; the sum is non-negative in both components.
\end{proof}


\section{Proposition 11.3 — tail symbol characterisations}

\begin{lemma}[Prop 11.3 — characterisation of $x_\tau$, $x_1$ via signatures and $\eps$]
  \label{lem:11_3}
  \lean{epsilon_zero_iff_tailSymbol_d, tailSig_fst_zero_iff_tailSymbol_s, tailSig_snd_zero_iff_tail_all_r}
  \leanok
  \uses{def:tailSymbol, def:tailSignature, def:epsilon_pbp, def:tailContrib}
  Assume $\mathcal{P}.\mathrm{colLen}(0) > \mathcal{Q}.\mathrm{colLen}(0)$. Then:
  \begin{enumerate}
    \item $\eps_\tau = 0 \iff x_\tau = d$.
    \item $p_{\tau_t} = 0 \iff x_\tau = s \iff $ every tail cell is $s$.
    \item $q_{\tau_t} = 0 \iff $ every tail cell is $r$ $\iff x_1 = r$ and $x_\tau = r$.
  \end{enumerate}
\end{lemma}

\begin{proof}
  \leanok
  \uses{def:tailContrib}

  \textbf{(1)} By layer monotonicity in column~0, a symbol $d$ (with maximal
  $\mathrm{layerOrd} = 4$) can only appear at the bottom.
  Conversely, $\eps_\tau = 0$ iff $d$ is present somewhere in column~0, which
  by monotonicity forces it to the bottom, i.e., $x_\tau = d$.

  \textbf{(2)} $p_{\tau_t} = \sum_i \mathrm{tailContrib}(x_i).p$, and
  $\mathrm{tailContrib}(\sigma).p = 0$ only for $\sigma = s$. Combined with
  the non-negativity of each term, $p_{\tau_t} = 0 \iff $ every tail cell is $s$.

  Now the equivalence ``every tail cell is $s$'' $\iff$ $x_\tau = s$ follows
  from layer monotonicity: $x_\tau = s$ means the $\mathrm{layerOrd}$ at the
  bottom is $1$; all cells above have $\mathrm{layerOrd} \leq 1$; combined
  with non-dot (every tail cell is non-$\bullet$ by $\mathrm{dot\_match}$),
  every cell has $\mathrm{layerOrd} = 1$, i.e., equals $s$.

  \textbf{(3)} $q_{\tau_t} = \sum_i \mathrm{tailContrib}(x_i).q$, and
  $\mathrm{tailContrib}(\sigma).q = 0$ only for $\sigma = r$. So
  $q_{\tau_t} = 0 \iff$ every tail cell is $r$.

  The second equivalence: ``every tail cell is $r$'' $\iff$ $x_1 = r$ and
  $x_\tau = r$. Layer monotonicity gives
  $\mathrm{layerOrd}(x_1) \leq \mathrm{layerOrd}(x_i) \leq \mathrm{layerOrd}(x_\tau)$
  for every tail cell $x_i$. If $x_1 = x_\tau = r$ (both equal $2$), then
  every $x_i$ has $\mathrm{layerOrd}$ sandwiched as $2 \leq \mathrm{layerOrd}(x_i) \leq 2$,
  so $\mathrm{layerOrd}(x_i) = 2$, i.e., $x_i = r$. The converse is trivial.
\end{proof}

\begin{remark}
  Unlike part~(2), part~(3) cannot be stated purely in terms of the bottom
  symbol $x_\tau$: the tail can have an $s$ cell above an $r$ bottom (layer
  orders $1 \leq 2$), giving $q_{\tau_t} = 2 \neq 0$. The correct statement
  requires either every tail cell to be $r$ or, equivalently, both $x_1$ and
  $x_\tau$ to be $r$.
\end{remark}

\section{Signature sum and residual identity}

\begin{lemma}[D-type signature sum]
  \label{lem:sig_sum_D}
  \lean{PBP.signature_sum_D}
  \leanok
  \uses{def:signature_pbp, def:cardCells}
  For $\tau \in \PBP_D$:
  \[ p_\tau + q_\tau = 2\cdot|\tau|. \]
  (Each cell contributes exactly $2$ to the total $p + q$.)
\end{lemma}
\begin{proof}
  \leanok
  Combine $p_\tau = \#_\bullet + \#_s + 2\#_r + \#_c + \#_d$ and the symmetric expression for $q_\tau$. Since $\mathcal{Q}$ is all $\bullet$ for D type, the non-dot counts come from $\mathcal{P}$ only, and $p_\tau + q_\tau = 2 \cdot |\mathcal{P}| + 2 \cdot |\mathcal{Q}| = 2|\tau|$.
\end{proof}

\begin{lemma}[Tail residual identity]
  \label{lem:residual_D}
  \lean{PBP.residual_identity_D}
  \leanok
  \uses{def:tailSignature, lem:tailSig_sum}
  For $\tau \in \PBP_D$:
  \[ p_{\tau_t} + q_{\tau_t} = 2 \cdot \mathbf{c}_1(\mathcal{O}) - 2 \cdot \mathbf{c}_2(\mathcal{O}). \]
\end{lemma}
\begin{proof}
  \leanok
  $k = \mathbf{c}_1 - \mathbf{c}_2$ by the BV dual column-length relations (specifically $\mathcal{P}.\mathrm{colLen}(0) = \mathbf{c}_1$, $\mathcal{Q}.\mathrm{colLen}(0) = \mathbf{c}_2$). Apply Lem.~\ref{lem:tailSig_sum}.
\end{proof}

\section{Proposition 11.4 — signature decomposition}

\begin{proposition}[Prop 11.4 D — signature decomposition (Eq 11.7)]
  \label{prop:11_4}
  \lean{PBP.signature_fst_decomp_D, PBP.signature_snd_decomp_D, prop_11_4_signature_decomp_D}
  \leanok
  \uses{def:signature_pbp, def:tailSignature, def:doubleDescent_D}
  (Primitive case $\mathbf{r}_2(\check{\mathcal{O}}) > \mathbf{r}_3(\check{\mathcal{O}})$.) The signature splits as
  \[ \Sign(\tau) = (\mathbf{c}_2(\mathcal{O}), \mathbf{c}_2(\mathcal{O})) + \Sign(\nabla^2\tau) + \mathrm{tailSig}(\tau). \]
  More explicitly (Lean form):
  \begin{align*}
    p_\tau &= \mathbf{c}_2(\mathcal{O}) + p_{\nabla^2\tau} + p_{\tau_t} \\
    q_\tau &= \mathbf{c}_2(\mathcal{O}) + q_{\nabla^2\tau} + q_{\tau_t}
  \end{align*}
  where $\mathbf{c}_2(\mathcal{O}) = \tau.\mathcal{Q}.\mathrm{colLen}(0)$.
\end{proposition}

\begin{proof}
  \leanok
  Decompose $\mathcal{P}$'s cells into three groups:
  \begin{enumerate}
    \item \emph{Column~0 below $\mathcal{Q}$'s col~0} ($i < \mathcal{Q}.\mathrm{colLen}(0)$): these are all $\bullet$ by the dot\_match axiom, contributing $\mathbf{c}_2(\mathcal{O}) \cdot (1, 1) / 2$ per cell... actually the direct expansion gives $\mathbf{c}_2(\mathcal{O})$ to both $p$ and $q$ via the $\bullet$ contribution to the signature formula.
    \item \emph{Columns $\geq 1$}: these match $\nabla^2\tau$'s columns exactly (up to shift), contributing $\Sign(\nabla^2\tau)$ (on a shifted shape).
    \item \emph{Tail cells} ($\mathcal{Q}.\mathrm{colLen}(0) \leq i < \mathcal{P}.\mathrm{colLen}(0)$ in column~0): contribute $\mathrm{tailSig}(\tau)$ by Def.~\ref{def:tailSignature}.
  \end{enumerate}
  Uses \texttt{countSym\_split}, \texttt{col0\_dot\_below\_Q\_D}, \texttt{Q\_countSym\_eq\_zero\_of\_D}, and the per-component $\mathrm{tailContrib}$ table.
\end{proof}
